\section{Limitations}
\label{sec:limitations}
% Each limitation is stated with its consequence for how the results must be read (SAP_v2.md, SAP_v3.md, SAP_v4.md, deviations log).
\begin{enumerate}
\item \emph{Synthetic supervision and one source of real pairs.} The bone image of a real radiograph is never observed. Real-image fidelity
 is measured against dual-energy soft-tissue images from one hospital and one acquisition set-up (GE Discovery XR656), which differ from the
 radiograph in tone and fine texture, so raw and band-limited metrics are both reported and calibrated metrics are not absolute.
\item \emph{The audit is linear and additive, and shares its mechanism with the regulariser.} Lesions are added in image units; beam
 hardening, scatter and tone mapping are not modelled, so $A$ measures the leak of an added signal rather than clinical detectability, and it
 is read next to a suppression gain because $A=0$ is reached by doing nothing. The real-nodule contrast and downstream subgroup analyses
 limit, but do not remove, the circularity of inserting lesions both in training and in the audit.
\item \emph{Development use of SZCH.} All 75 SZCH images were used in development (audit) and 24 for paired decisions, so only the 51
 held-out pairs support fidelity claims; the plan changed once after the SZCH pairs had been looked at.
\item \emph{Resolution, backbone and operating point.} Our suppressors predict at 512\,px with tiny encoders, the reference model at
 1024\,px with SwinV2-B. The SSIM gap sits in fine texture (with a band limit of $\sigma=4$ it shrinks from 0.021 to 0.002); whether
 prediction at 1024\,px closes it was not tested. Only $w=16$ was evaluated with the full protocol.
\item \emph{Uncertainty.} Conformal calibration treats pixels as exchangeable and was calibrated on synthetic data; the recalibration on real
 pairs was chosen after the shift was seen and used 24 pairs from one hospital. Ensemble members share training data.
\item \emph{Scope and external validity.} Training and paired evaluation use adult frontal radiographs; BDCXR-3257 is a single-centre
 paediatric set with image-level labels and no patient identifiers, and its audit inserts adult nodules into children's radiographs. There
 is no lateral imaging and no reader study. On JSRT the bone overlap of each audit site came from the reference model's bone map. Downstream
 models were not tuned per input, and the overlap analysis is not estimable on TBX11K. The models are for research use only.
\item \emph{Classifier study.} One encoder (DenseNet-121) and one resolution ($256\times256$\,px) were tested; heavier encoders or full
 resolution may use the extra streams differently. The NIH test split served the downstream experiments of Study I as well, its labels are
 report-derived, and VinBigData has one focal class. The per-class DeLong tests treat NIH images as independent although patients contribute
 several images, and the pointing game depends on the attribution method, the 15\,px tolerance and the eight annotated NIH classes.
\item \emph{Cross-modality study.} The 2.5D CT recipe was seed-unstable, and the replacement of seed 2 by seed 5 followed a rule fixed
 beforehand but used a teacher trained later. The CT subset is small (\num{1600} training volumes) and its labels are report-derived. The
 localisation of the H13 difference in hernia (86 positive test images) comes from per-class tests planned only as descriptive, and H13
 compares encoders trained with different pre-specified budgets (DenseNet-121 for 20 epochs at the full encoder learning rate, Swin for 10
 epochs at a quarter of it), neither tuned further; the reference's NIH test score was known from Study II when the fourth plan was frozen.
 H14 relies on
 Grad-CAM++ for both modalities, the least localising of the three maps on real radiographs. The example heat maps contain no mass or nodule,
 and VinBigData fine-tuning was not run. Teacher outputs, predictions and weights reached the authors' computer against the design's own
 data rule; nothing derived from CT-RATE was released.
\item \emph{Post-freeze corrections and registration.} Analysis errors found in an independent check after the first locked run were
 corrected (among them a resampling error that narrowed interaction intervals and detection metrics that depended on the order of tied
 scores by up to 0.00055); cloud and local GPU outputs differ by at most 0.0024\,dB. No pre-specified verdict changed. The analysis plans were
 hash-frozen in version control before testing but not deposited in a public registry.
\end{enumerate}
